\section{Resource Burdens of NAS in FL}

NAS may train or estimate the performance of many candidates, while distributing this work across clients adds resource costs to an already expensive procedure~\cite{ying2019nasbench101reproducibleneuralarchitecture,fl_advances_and_open_problems_2021}. This theme covers challenges that make federated architecture search operationally burdensome and solutions that reduce, reuse or reorganise the required work.

\subsection{Computation and Search Effort}

\subsubsection{Challenge 1.1: Client Resource Limits Make Search Workloads Infeasible}

NAS-in-FL becomes infeasible when an assigned search workload exceeds a client's computation, memory or per-round time budget. Absolute resource limits are compounded by hardware heterogeneity because a workload that is feasible for one client may exceed another client's capacity~\cite{fedoras_2022,fednas_over_hetero_mobile_devices_2022,dfnas_2020,peaches_2024,supernet_trade-off_fednas_2025}. The problem is particularly acute for differentiable supernet-based strategies that distribute a shared supernet and require clients to update several mixed candidate operations locally~\cite{fednas_2021,dp-fnas_2020,feathers_2023,fed_auto_mri_2023}. Other strategies usually assign one discrete or generated candidate at a time, but even one candidate can exceed the resources available to a weak client. This claim concerns whether the assigned work is feasible. Differences in completion time among clients that can perform the work are treated in Challenge 4.2.

\noindent\textbf{Solution 1.1a: Build the search space from inexpensive components.} Practitioners can apply this solution to every search strategy type by restricting candidate architectures to inexpensive operations, blocks or modules, depending on the search space type. Prominent examples include depthwise-separable convolutions and mobile-oriented CNN blocks~\cite{fed_auto_mri_2023,intrusion_detection_fednas_2024}. This makes client-side evaluation cheaper, but may exclude expensive components that increase prediction accuracy.

\noindent\textbf{Solution 1.1b: Match sampled search work to client capacity.} Supernet-based strategies can assign each client a feasible subspace or sampled subnet instead of a weighted mixture of all candidate operations. The local depth, width or optimisation budget can be adjusted to the client's computation and memory~\cite{dfnas_2020,fedoras_2022,apons_2023,divide-and-conquer_fednas_2023,peaches_2024,fgfl_2024,supernet_trade-off_fednas_2025,optimize_fednas_edge_2021,efnas_2024,load_monitoring_fednas_2025}. This allows weak and strong clients to contribute without executing the complete supernet. However, sampling only part of the supernet can create unequal operation and width coverage, which is addressed separately in Challenge 3.2.

\noindent\textbf{Solution 1.1c: Search trainable adaptation modules over a frozen backbone.} When a suitable pretrained model is available, clients can keep its backbone fixed and search a sparse combination of lightweight adaptation modules. Only the selected modules and their architecture variables require local optimisation and exchange, reducing both client computation and update size~\cite{fednaspet_2023}. The benefit depends on how well the fixed backbone transfers to the clients' tasks.

\noindent\textbf{Solution 1.1d: Keep the full search model on the server.} The server can retain the complete supernet and send each client only a feasible sampled subnet~\cite{perfedrlnas_2024,rafl_2024}. Clients can also train lightweight knowledge models whose predictions guide server-side supernet training through distillation~\cite{hapfnas_2025}. These designs keep the largest search structure off constrained devices.

\subsubsection{Challenge 1.2: Repeated Federated Candidate Evaluation Makes Search Costly}

Even when each local evaluation is feasible, comparing architectural alternatives across repeated federated rounds consumes substantial computation and elapsed time. Black-box strategies may train several discrete candidates through separate federated evaluations~\cite{decnas_2022,rt_fed_evo_nas_2020}, while weight-sharing strategies repeatedly update an over-parameterised supernet or its sampled subnets~\cite{fednas_2021,superfednas_2024}. Sequential split-model execution can also leave clients idle while other model segments run~\cite{nasfly_2024}. Energy provides an additional measure of this work because local optimisation consumes computation energy and repeated exchanges consume transmission energy~\cite{f2nas_2024}. Communication volume itself is treated separately in Challenge 1.5. A distinct post-search training cost applies when the selected architecture cannot reuse search-time weights, as discussed in Challenge 1.3.

\noindent\textbf{Solution 1.2a: Dynamic candidate evaluation budgets.} This solution applies to black-box search strategies that compare discrete candidates through independent training, including greedy pruning, evolutionary and genetic algorithms, reinforcement learning and Bayesian optimisation. Similar to multi-armed-bandit AutoML approaches ~\cite{automl_book_hpo_2019}, each candidate initially receives only a small budget of federated rounds, after which weak candidates can be discarded and promising candidates trained for longer.

For example, \cite{decnas_2022,fednas_over_hetero_mobile_devices_2022} observes that the best candidate usually becomes distinguishable within the first couple of communication rounds, drops the remaining candidates and increases the round budget in later search iterations. Evolutionary searches similarly use a fixed small number of rounds or early stopping to limit candidate evaluation~\cite{multi_objective_evo_fl_2020,gpt4-boosted_pso_fednas_2024}.

\noindent\textbf{Solution 1.2b: Prune the supernet progressively.} Supernet-based strategies can periodically remove weak candidate operations from the shared supernet, so fewer alternatives remain stored and trained in subsequent federated rounds~\cite{dfnas_2020,peaches_2024}. Unlike dynamic candidate evaluation budgets, this changes the search space itself rather than the number of rounds allocated to independently trained candidates. Its savings increase as the search progresses, but the initial rounds still incur the cost of the complete supernet and early pruning can permanently remove useful operations.

\noindent\textbf{Solution 1.2c: Predict candidate performance.} Performance prediction is a standard AutoML technique for reducing the number of costly configuration evaluations~\cite{automl_book_hpo_2019}. In NAS-in-FL, candidate-based strategies can evaluate a small set of architectures and train a predictor to estimate the quality of the remaining candidates~\cite{hapfnas_2025,energy_fednas_2025}. Additional care is required because private and non-IID client data can make architecture rankings client-specific. Predictors can therefore be trained separately for each client~\cite{hapfnas_2025}, or learned collaboratively by aggregating their parameters without transmitting the locally evaluated architectures~\cite{federated_bayesian_optimization_nas_2023}. This reduces the number of candidates evaluated by clients, but its reliability depends on whether the evaluated architectures represent the relevant parts of the search space.

\noindent\textbf{Solution 1.2d: Evaluate candidates in parallel.} Since black-box search strategies generate several distinct candidates per iteration, the candidates can be trained and evaluated concurrently by separate client groups~\cite{decnas_2022,rt_fed_evo_nas_2020,energy_fednas_2025}.

Supernet-based strategies can instead sample several subnets, allocate each to a client cluster and merge their updates back into the shared supernet~\cite{finch_2024}. This potentially reduces elapsed search time when enough clients are available, but comparisons can become biased when the groups evaluate candidates on different data distributions~\cite{decnas_2022}.

\noindent\textbf{Solution 1.2e: Reuse supernet weights for candidate evaluation.} Weight-sharing strategies can train common operation weights and allow each candidate to inherit the subset associated with its architecture~\cite{fl-agnns_2023,optimize_fednas_edge_2021}. This avoids training every candidate independently from scratch, although inherited-weight performance can rank candidates differently from full training. A short incremental federated fine-tuning stage can provide a higher-fidelity estimate for a smaller set of promising candidates~\cite{medical_image_segmentation_fednas_2024}.

\noindent\textbf{Solution 1.2f: Supervise concurrent search agents.} Several search agents can run in parallel while a supervisor compares their relative losses and adjusts the momentum applied by each agent. The final architecture can then be selected according to prediction error, training time and model size~\cite{supervised_fednas_2023}. This can smooth and accelerate convergence when parallel resources are available, but requires maintaining several search trajectories.

\noindent\textbf{Solution 1.2g: Use split-training idle periods for local search.} Auxiliary input and output components can turn an assigned model segment into a locally executable supernet. During otherwise idle periods, a client can sample and train alternative branches while using the corresponding segment of the split model as a teacher for feature distillation~\cite{nasfly_2024}. This converts waiting time into search work that can proceed without the remaining model partitions.

\noindent\textbf{Solution 1.2h: Adapt local epochs to an energy--convergence bound.} The number of local epochs can be selected in each round from an upper bound that relates loss reduction to client computation and communication costs. This balances additional local work against another cloud--client exchange rather than applying one fixed value throughout the search~\cite{f2nas_2024}. The cited evaluation reports an energy reduction of nearly 50\%, although applying the mechanism requires estimates of client energy and latency.

\subsubsection{Challenge 1.3: Post-Search Retraining Duplicates Federated Computation}

Some NAS strategies use search-time training only to select an architecture and then train the selected model from a new initialisation. In NAS-in-FL, this adds another sequence of client-side epochs and aggregation rounds after clients have already spent computation and communication on the search~\cite{medical_image_segmentation_fednas_2024,privacy_preserving_fednas_for_edge_2025}. This challenge applies whenever the search produces an architecture without reusable final weights and therefore delays the availability of a deployable model.

\noindent\textbf{Solution 1.3a: Reuse weights learned during search.} Weight-sharing supernet strategies can initialise the selected subnet with the corresponding operation weights learned during search and perform only limited final fine-tuning~\cite{medical_image_segmentation_fednas_2024}. This reuses client computation that has already been spent on the selected operations, although weights learned jointly with competing subnets may be a poor initialisation for the final model.

\noindent\textbf{Solution 1.3b: Produce a deployable model during search.} One-stage supernet-based strategies can jointly optimise architecture selection and model weights while concentrating the architecture distribution or pruning weak operations, so the search terminates with a discrete model and trained weights~\cite{dfnas_2020,fed_meta_nas_2025}. Reinforcement-learning strategies over a weight-sharing supernet can similarly turn later search rounds into fine-tuning as their sampling probabilities concentrate~\cite{perfedrlnas_2024}. This eliminates the separate retraining stage, but the final path may remain undertrained if it received too few updates while sharing weights with competing paths.

\subsubsection{Challenge 1.4: Heterogeneous Deployment Requirements Multiply Search Cost}

A NAS-in-FL deployment may include devices with different inference-time computation, memory or latency constraints. One architecture sized for the weakest client can underuse stronger hardware, while one sized for the strongest client can be infeasible elsewhere~\cite{fedoras_2022}. The search must therefore produce several target-specific outcomes. Obtaining each through an independent federated search repeats client computation and communication, so the total cost grows with the number of device classes or deployment budgets~\cite{superfednas_2024}.

\noindent\textbf{Solution 1.4a: Search architectures for hardware tiers within shared training.} Clients with similar inference constraints can be assigned to capability tiers, with one feasible architecture searched for each tier~\cite{fedoras_2022,clustered_direct_fednas_2022}. Sharing supernet weights across the tier-specific searches prevents each outcome from requiring an independent federated training process.

\noindent\textbf{Solution 1.4b: Train one supernet for constraint-specific selection.} A weight-sharing supernet can expose subnets with different depths, widths or exits. After federated training, each client can select a subnet that satisfies its hardware limits and performs well on its validation data~\cite{superfednas_2024,hapfnas_2025,enasfl_2024}. Predictor-guided local selection can avoid further federated training and has reduced training cost elevenfold for 20 targets~\cite{superfednas_2024}. Co-training many subnets can, however, cause weight-sharing interference and leave rarely sampled parameters stale~\cite{superfednas_2024}.

\noindent\textbf{Solution 1.4c: Include target-specific hardware costs in the search objective.} Candidate-based and supernet-based strategies can optimise predictive performance subject to a measured latency or multiply--accumulate operation budget~\cite{clustered_direct_fednas_2022,load_monitoring_fednas_2025,fgfl_2024}. Supplying a hardware profile or budget for each client group makes the selected outcomes specific to their deployment constraints within one coordinated search.

\subsection{Communication Burden}

Cross-device clients often communicate with the server over capacity-limited wide-area links. Typical wide-area bandwidths of 5--25 Mbit/s are far below the more than 10 Gbit/s available within data centres~\cite{peaches_2024}. The discussion below concerns the communication added by search and by the selected model. Differences between client links are treated under disrupted architecture search, where they affect timeliness, participation and candidate assignment.

\subsubsection{Challenge 1.5: Architecture Search Multiplies Communication Load}

Fixed-model FL exchanges the parameters of one architecture. NAS-in-FL can add both larger payloads and more frequent synchronisation. A supernet contains weights for mutually exclusive operations and can be much larger than any selected architecture, while differentiable strategies may exchange architecture variables alongside model weights throughout the search~\cite{fedoras_2022,hapfnas_2025,fednas_2021,fedregnas_2025}. The resulting communication load can exceed client transfer budgets even when each individual architecture is feasible.

\noindent\textbf{Solution 1.5a: Compress search updates.} Quantisation represents update entries with fewer bits, while sparsification transmits only selected entries. One differentiable supernet-based approach uses 8-bit stochastic quantisation and retains the largest 20\% of entries after clipping~\cite{fedregnas_2025}. These operations reduce the transferred weight and architecture updates without changing the search variables.

\noindent\textbf{Solution 1.5b: Encode architecture choices compactly.} When clients already hold a compatible master model or trained supernet, the server can identify a candidate using a selector or genotype instead of transmitting its parameters. Clients reconstruct the indicated subnet locally and return its evaluation or updated weights~\cite{rt_fed_evo_nas_2020,medical_image_segmentation_fednas_2024,intrusion_detection_fednas_2024}. This requires the server and clients to share the search-space definition and clients to retain the common weights.

\noindent\textbf{Solution 1.5c: Synchronise architecture variables less frequently.} Architecture variables can be treated as slow variables that are updated and uploaded less frequently than model weights. One approach transmits them every \(H_{\alpha}\) rounds and stops doing so after the architecture has been fixed, while ordinary weight training continues~\cite{fedregnas_2025}.

\noindent\textbf{Solution 1.5d: Aggregate search updates hierarchically.} A nearby aggregator can combine the model and architecture updates of several clients before forwarding one result to the remote server. Frequent aggregation then uses local links, while fewer exchanges traverse the wide-area network~\cite{finch_2024}. In one non-IID CIFAR-10 evaluation, this reduced traffic at 80\% target accuracy from 3.29--4.68 GB to 1.64 GB~\cite{finch_2024}.

\subsubsection{Challenge 1.6: Selected Architectures Can Exceed Federated Training Transfer Budgets}

An architecture selected only for predictive performance or local computation can still contain many parameters that FL must transmit in every subsequent training round. The search decision therefore determines the continuing load on client links, and a model that meets computation and memory budgets may still exceed the transfer budget~\cite{multi_objective_evo_fl_2020}.

\noindent\textbf{Solution 1.6a: Optimise for transmitted model size.} Communication cost can be included directly in the architecture objective. Evolutionary search can optimise global error together with the average number of weights uploaded by clients, favouring sparse architectures~\cite{multi_objective_evo_fl_2020}. Differentiable supernet-based search can similarly regularise the expected number of transmitted bytes~\cite{fedregnas_2025}. Both expose the trade-off between predictive performance and continuing communication cost before selecting an architecture.
